\section{Connectivness of Graphs}

\begin{definition}
Vertices $u, v \in V, u \neq v$ are called \textbf{connected} if there is a $u-v$-path in $G$.
\end{definition}

\begin{definition}
A graph $G = (V, E)$ is \textbf{connected} if $\forall u, v \in V, u \neq v$ a $u-v$-path in $G$ exists.
\end{definition}

\begin{definition}
Vertices $u, v \in V, u \neq v$ are called \textbf{$k$-connected} if $|V|\ge k+1$ and
$$\forall X \subseteq V \setminus \{u, v\}$$
$$ \text{ with } |X| \lt k,$$
$$ u \text{ and } v \text{ are connected in } G[V \setminus X]$$
\end{definition}

\begin{definition}
A graph $G = (V, E)$ is called \textbf{$k$-connected} if $|V| \gt k$ and
$$\forall X \subseteq V \text{ with } |X| \lt k$$
$$ G[V \setminus X] \text{ is connected}$$
\end{definition}

\begin{definition}
$X \subseteq V \setminus \{u, v\}$ is called a \textbf{$u-v$-separator} if $u$ and $v$ are in different connected components of $G[V \setminus X]$.
\end{definition}

\begin{remark}
If $G$ contains a vertex $v$ with $deg(v) < k$, then $G$ cannot be $k$-connected.
\end{remark}

\begin{theorem}
\textbf{Menger's Theorem (Vertex Version)}: The minimum size of a $u-v$-separator is equal to the maximum number of internally vertex-disjoint $u-v$-paths. A graph is $k$-connected iff $\forall u \neq v$ there are $k$ internally vertex-disjoint $u-v$-paths.
\end{theorem}

\begin{definition}
A graph $G = (V, E)$ is called \textbf{$k$-edge-connected} if
$$\forall X \subseteq E \text{ with } |X| \lt k$$
$$(V, E \setminus X) \text{ is connected}$$
\end{definition}

\begin{theorem}
\textbf{Menger's Theorem (Edge Version)}: A graph $G$ is $k$-edge-connected iff. $\forall u, v \in V, u \neq v$ there are $k$ edge-disjoint $u-v$-paths.
\end{theorem}

\begin{theorem}
It always holds that:
$$\text{Vertex Connectivity} \le \text{Edge Connectivity}$$
$$\le \text{Minimum Degree } \delta(G)$$
\end{theorem}

Lets look at the special case where $k=1$:

\begin{definition}
For a connected graph $G = (V, E)$:
\begin{itemize}
\item Vertex $v \in V$ is a \textbf{cut vertex} iff. $G[V \setminus \{v\}]$ is disconnected.
\item Edge $e \in E$ is a \textbf{bridge} iff. $(V, E \setminus \{e\})$ is disconnected.
\end{itemize}
\end{definition}

\begin{lemma}
If $e = \{x, y\} \in E$ is a bridge, then $\text{deg}(x) = 1$ or $x$ is a cut vertex.
\end{lemma}

\begin{definition}
\textbf{Blocks}: Define an equivalence relation $\sim$ on $E$:
$$e \sim f \iff e=f \lor \exists \text{ cycle containing } e \text{ and } f$$
The equivalence classes are called \textbf{blocks}.
\end{definition}

Note: Two blocks can share at most one vertex, which must be a cut vertex.

\begin{definition}
\textbf{Block Graph}: Bipartite graph $T = (A \cup B, E_T)$ where:
$A = \{\text{cut vertices}\}$
$B = \{\text{blocks}\}$
Edge $\{a, b\} \in E_T \iff a$ is incident to an edge in block $b$.
\end{definition}

\begin{theorem}
If $G$ is connected, the block graph of $G$ is a tree.
\end{theorem}

\begin{algorithm}
#DFS & Low Values for cut vertices/bridges in O(|V|+|E|)
low[v] := #smallest DFS number reachable from v by tree edges and <=1 back edge.

low[v] = min{ dfs[v], min_{(v,w) back} dfs[w], min_{(v,w) tree} low[w] }

#A vertex v is a cut vertex iff:
v != root and exists child u with low[u] >= dfs[v].
v = root and has >= 2 children in DFS tree.

#A tree edge (v,w) is a bridge iff:
low[w] >= dfs[v].
\end{algorithm}
